\documentclass[11pt,a4paper]{article}

\usepackage[top=2cm, bottom=2cm, left=2cm, right=2cm]{geometry}
\usepackage[T1]{fontenc}
\usepackage{mathpazo}
\usepackage[scaled=0.92]{helvet}
\usepackage{enumitem}
\usepackage{hyperref}
\usepackage{parskip}
\usepackage{microtype}
\usepackage{float}
\usepackage[most]{tcolorbox}

\definecolor{accent}{HTML}{7DBBB5}
\definecolor{accentdark}{HTML}{3D5A5B}
\definecolor{accentlight}{HTML}{D4E8E5}
\definecolor{faqq}{HTML}{FDF8F3}
\definecolor{faqborder}{HTML}{7DBBB5}
\definecolor{sectionrule}{HTML}{EDE5DD}
\definecolor{titlebg}{HTML}{3D5A5B}
\definecolor{knownbg}{HTML}{FDF8F3}
\definecolor{knownborder}{HTML}{E8836B}
\definecolor{lockbg}{HTML}{D4E8E5}
\definecolor{lockborder}{HTML}{7DBBB5}
\definecolor{bodytext}{HTML}{5C6B6B}
\definecolor{coral}{HTML}{E8836B}
\definecolor{teal}{HTML}{7DBBB5}
\definecolor{purple}{HTML}{9B8EC4}

\hypersetup{
    colorlinks=true,
    linkcolor=accentdark,
    urlcolor=accent,
    hypertexnames=false
}

\setlist[itemize]{itemsep=3pt, parsep=1pt, leftmargin=1.6em}

% Subsection heading with a screenshot sitting directly next to the heading text.
% #1 = heading text, #2 = image path
\newcommand{\imgsubsection}[2]{%
    \subsection*{#1\hspace{3em}%
        \raisebox{-0.3\height}{\includegraphics[height=0.7cm]{#2}}%
    }%
}

\newtcolorbox{titlebox}{
    colback=faqq,
    colframe=faqborder,
    coltext=accentdark,
    boxrule=3pt,
    arc=4pt,
    left=10pt, right=10pt, top=8pt, bottom=8pt
}

\newtcolorbox{faqbox}[1]{
    colback=faqq,
    colframe=faqborder,
    boxrule=0.6pt,
    arc=3pt,
    left=8pt, right=8pt, top=6pt, bottom=6pt,
    fonttitle=\bfseries,
    coltitle=accentdark,
    toptitle=3pt, bottomtitle=2pt,
    colbacktitle=accentlight,
    title={#1}
}

\newtcolorbox{knownbox}{
    colback=knownbg,
    colframe=knownborder,
    boxrule=0.6pt,
    arc=3pt,
    left=8pt, right=8pt, top=6pt, bottom=6pt
}


\newcommand{\sectionruleline}{%
    \vspace{2pt}%
    {\color{sectionrule}\hrule height 0.5pt}%
    \vspace{6pt}%
}

\pagestyle{plain}

\begin{document}

\begin{titlebox}
    \centering
    {\LARGE\textsf{\textbf{Sexual Health App User Guide}}}\\[0.5em]
    {\normalsize SDDP Group 30\enspace$\cdot$\enspace COMP2300\enspace$\cdot$\enspace University of Southampton}
\end{titlebox}

\vspace{8pt}

\section*{Part 1 - Getting Started}
\sectionruleline

\subsection*{First Launch and Setup}

When you open the app for the first time, a short onboarding walkthrough introduces
core features and personalisation.
After the walkthrough you are asked to create a \textbf{secret equation},
a simple calculation you will remember to unlock the app, such as \texttt{7\,+\,6\,=\,13}.
Type the equation on the built-in calculator keypad, press~\texttt{=},
review the confirmation prompt, and press \emph{Confirm}.

From this point on, the app always opens as a calculator.
To access the real app, enter your secret equation and press~\texttt{=};
the main interface is revealed with a smooth transition.

\imgsubsection{Article Tab}{D7//figures/articlemenu.png}


\vspace{10pt}

The \textbf{Articles} tab is the default view after unlocking.

\begin{itemize}
    \item A personalised \emph{Articles For You} feed is shown on the home screen,
          ranked by relevance to your preferences, and your reading history.
    \item Use the \textbf{search bar} at the top to find articles by keyword
          (e.g.\ ``contraception'', ``STI testing'').
          Results are ranked using a hybrid of text matching and meaning-based search,
          so everyday language works well.
    \item Tap an article to open it. Each article is split into
          \textbf{swipeable sections} so you are never faced with a wall of text.
          A \textbf{section navigation menu} at the top lets you jump directly to
          any section heading.
    \item \textbf{Related articles} are suggested at the bottom of each article page.
    \item A \textbf{Recently Read} area on the article list screen shows articles you have
          previously opened. Tapping one resumes from where you left off.
\end{itemize}


\imgsubsection{Calendar Tab}{D7//figures/calmenu.png}

\vspace{10pt}

The \textbf{Calendar} tab lets you manage appointments, medication, and tests.

\begin{itemize}
    \item A \textbf{monthly grid} shows the current month.
          Days that contain events display colour-coded indicators:
          {\color{coral}\textbf{coral}} circles for Appointments, {\color{teal}\textbf{teal}} triangles for Medication,
          and {\color{purple}\textbf{purple}} squares for Tests. The different shapes ensure accessibility
          for colour-blind users.
    \item Tap a day cell to see that day's events listed below the grid.
          Use the arrows or the \textbf{month/year picker} to navigate between months.
    \item Tap the \textbf{+ new event}~button to create a new event.
          Fill in a name, date, and event type.
          Optional fields include time, description, dosage (for medications),
          and a recurrence rule (daily, weekly, monthly, or yearly).
    \item Recurring events appear automatically on every matching date in the
          calendar and event feed.
    \item The \textbf{Event Feed} (accessible from the calendar) lists all upcoming
          events in chronological order. Tap any event to open its
          \textbf{detail page}, where you can view full information, edit, or delete it.
    \item \textbf{Reminders} send discreet in-app notifications at
          scheduled times. Notification wording is kept vague by default so it does
          not reveal the app's purpose.
\end{itemize}

\imgsubsection{Settings Tab}{D7//figures/setmenu.png}

\vspace{10pt}

The \textbf{Settings} tab lets you personalise the app.

\subsubsection*{Content Preferences}
Open \emph{Content preferences} to control which articles appear across browse,
search, and recommendations.
\begin{itemize}
    \item \textbf{Blocked tags:} select tags for topics you want to hide.
          Articles matching any blocked tag are removed from all article lists,
          search results, and recommendations.
    \item \textbf{Prioritised tags:} select tags for topics that matter to you.
          Matching articles receive a ranking boost in search and a stronger
          result highlight.
\end{itemize}
Both pickers offer a search field and tag suggestions. Blocking a tag
automatically removes it from the prioritised list, and vice versa.
\subsubsection*{Display}
Choose the visual style that suits you. A \emph{Reset to standard view} button
restores the default.
\begin{itemize}
    \item \textbf{Display mode:} switch between \emph{Standard}, \emph{Dark},
          and \emph{High Contrast} via radio buttons.
\end{itemize}
\subsubsection*{Text Size}
Adjust readability across the entire app. A \emph{Reset to defaults} button
restores both settings below.
\begin{itemize}
    \item \textbf{Global text size:} choose \emph{Small}, \emph{Standard},
          \emph{Large}, or \emph{Extra Large}.
    \item \textbf{Dyslexia-friendly font:} toggle the OpenDyslexic typeface for
          all readable content.
\end{itemize}
\subsubsection*{Reminders and Privacy}
\begin{itemize}
    \item \textbf{Off:} no pop-ups; events remain only in the calendar feed.
    \item \textbf{Disguised (Maximum Privacy):} notifications mimic system
          alerts. You can customise the disguise title and body text.
    \item \textbf{Discreet:} shows times only (e.g.\ ``Upcoming Event'').
    \item \textbf{Detailed:} shows the full event name and any notes.
\end{itemize}
\subsubsection*{Parental Controls}
Set a digits-only PIN to protect sensitive settings pages. Display and Text size
remain unlocked.
\subsubsection*{Privacy Policy}
A read-only page explaining local-only storage, zero analytics, on-device
notifications, and passcode encryption. Also contains a \textbf{Reset App}
action that permanently erases all data after a typed confirmation phrase.
\subsubsection*{App Disguise}
\begin{itemize}
    \item \textbf{Startup behaviour:} choose \emph{Enabled} (app opens to a
          calculator screen requiring your secret equation) or \emph{Disabled}
          (app opens directly to the home screen).
    \item \textbf{Lock-button behaviour:} choose whether the lock button returns
          to the calculator disguise or closes the application entirely.
\end{itemize}
\imgsubsection{Lock Tab}{D7//figures/lockmenu.png}

\vspace{10pt}

Tap the \textbf{lock icon} in the bottom navigation bar at any time to instantly
return to the calculator disguise. Useful for if someone walks in the room or you feel someone might snoop over your shoulder.

\vspace{10pt}
{\color{sectionrule}\hrule height 0.5pt}
\vspace{6pt}

\section*{Part 2 - FAQs and Known Issues}
\sectionruleline

\subsection*{Frequently Asked Questions}

\begin{faqbox}{Which operating systems can the app run on?}
Windows, macOS, and Linux.
The app requires \textbf{Java~17} or higher and \textbf{Maven~3.6+}.
Build with \texttt{mvn clean compile} and run with \texttt{mvn javafx:run}.
\end{faqbox}

\begin{faqbox}{What if I forget my secret equation?}
On the calculator screen, enter \texttt{999\,$\div$\,0} and press~\texttt{=}.
This resets your stored equation and returns you to the setup wizard so you can
create a new one.
\end{faqbox}

\begin{faqbox}{If I delete the app, does my data stay?}
Events are stored in local JSON files and your preferences use the operating
system's key-value store.
Removing the application may delete this data depending on your OS;
there is no cloud backup.
\end{faqbox}

\begin{faqbox}{Is my data sent anywhere?}
No. All data (events, content preferences, and your secret equation hash)
is stored locally on your device. The app does not use a network connection.
\end{faqbox}

\begin{faqbox}{Where does the article content come from?}
Articles are sourced from
\emph{Brook},
a reputable UK sexual health service. All content is packaged within the app and
loads without an internet connection.
\end{faqbox}

\begin{faqbox}{How do I get in touch if I encounter technical problems?}
You can contact us using our University of Southampton E-Mail addresses:\\
\texttt{jw14g24@soton.ac.uk},\enspace
\texttt{sh6n24@soton.ac.uk},\enspace
\texttt{tm2n24@soton.ac.uk},\enspace
\texttt{op2g24@soton.ac.uk}.
\end{faqbox}

\subsection*{Known Issues}

\begin{knownbox}
\begin{itemize}[topsep=0pt]
      \item Currently not buildable on iOS or Android
      \item On some text size / font sizes there may be some text truncation
\end{itemize}
\end{knownbox}

\end{document}